\documentclass[11pt]{article}
\usepackage[margin=1in]{geometry}
\usepackage{amsmath,amssymb,bm}
\usepackage{booktabs}
\usepackage{hyperref}
\usepackage{xcolor}
\usepackage{enumitem}

\title{Residual-Guided Hidden-Physics Discovery for PFM-Constrained BTO Phase-Field Simulations}
\author{sim\_bo / AutoPF--MatEnsemble campaign summary}
\date{June 22, 2026}

\newcommand{\vect}[1]{\bm{#1}}
\newcommand{\dd}{\mathrm{d}}
\newcommand{\eps}{\varepsilon}
\newcommand{\Pz}{P_z}
\newcommand{\Uz}{u_z}
\newcommand{\Uzc}[2]{u_{z,#2}^{#1}}

\begin{document}
\maketitle

\section{Purpose and Data Products}

The latest hidden-physics campaign was designed to answer a focused question:
after fixing the gradient coefficients to the dense-posterior, physically
tangible values
\[
  g_{11}=0.5,\qquad g_{12}=-0.06,\qquad g_{44}=0.02,
\]
what additional low-dimensional physics terms most consistently reduce the
PFM surface-displacement residuals?

The current best model is a co-active surface-screening field plus a localized
anisotropic, shear-like eigenstrain component.  The final full-holdout round is
\texttt{round\_anisotropic\_holdout\_full\_000}.  Its analysis files are under
\[
\texttt{sim\_bo\_bto\_physical\_bounds/anisotropic\_hidden\_analysis}.
\]
The MOOSE input used for the hidden-physics model is
\texttt{BTO\_DW\_anisotropic\_hidden\_physics\_zdecay\_ic.i}; the campaign
manifests were produced by
\texttt{build\_anisotropic\_hidden\_physics\_manifest.py} and
\texttt{build\_anisotropic\_holdout\_validation\_manifest.py}.

\section{Observation Model and Residual Objective}

For each experimental bias--pulse condition
\[
  c=(V,\tau),
\]
the experiment provides a surface PFM-derived observable, denoted
\(\Uzc{\mathrm{exp}}{c}(x,y)\).  The phase-field simulation gives a final
surface displacement \(\Uzc{\mathrm{sim}}{c}(x,y;\theta)\), where \(\theta\)
collects the gradient coefficients and hidden-physics parameters.

Before evaluating the objective, the simulated surface field is transposed into
the experimental orientation and resampled to the experimental grid.  Both
fields are then z-normalized condition-by-condition:
\[
  \widehat{u}_{z,c}^{\,\mathrm{sim}}
  =
  \frac{\Uzc{\mathrm{sim}}{c}-\langle \Uzc{\mathrm{sim}}{c}\rangle}
       {\mathrm{std}(\Uzc{\mathrm{sim}}{c})},
  \qquad
  \widehat{u}_{z,c}^{\,\mathrm{exp}}
  =
  \frac{\Uzc{\mathrm{exp}}{c}-\langle \Uzc{\mathrm{exp}}{c}\rangle}
       {\mathrm{std}(\Uzc{\mathrm{exp}}{c})} .
\]
The residual map is
\[
  R_c(x,y;\theta)
  =
  \widehat{u}_{z,c}^{\,\mathrm{sim}}(x,y;\theta)
  -
  \widehat{u}_{z,c}^{\,\mathrm{exp}}(x,y).
\]
The objective uses the center crop, excluding the outer 10\% edge band in each
image dimension.  If \(M(x,y)\in\{0,1\}\) is the crop mask, the normalized MSE is
\[
  \mathcal{L}(\theta)
  =
  \frac{1}{|\mathcal{C}|}
  \sum_{c\in\mathcal{C}}
  \frac{\sum_{x,y} M(x,y)\,R_c(x,y;\theta)^2}
       {\sum_{x,y} M(x,y)\,
       \left(\widehat{u}_{z,c}^{\,\mathrm{exp}}(x,y)\right)^2}.
\]
The edge masking was important because the largest residuals in the baseline
fits were often dominated by lateral/edge artifacts rather than by the domain
interior that should drive hidden-physics inference.

\section{Baseline Physics Model}

The core phase-field model is an LGD electro-elasto-static model with
polarization \(\vect{P}=(P_x,P_y,P_z)\), displacement \(\vect{u}\), and electric
potential \(\phi\).  Schematically, the relaxed state minimizes or evolves down
the functional
\[
  \mathcal{F}_0
  =
  \int_{\Omega}
  \left[
    f_{\mathrm{Landau}}(\vect{P})
    + f_{\mathrm{grad}}(\nabla \vect{P})
    + f_{\mathrm{elec}}(\vect{P},\phi)
    + f_{\mathrm{el}}(\vect{\eps},\vect{P})
  \right]\,\dd\Omega .
\]
The gradient term is controlled by the three cubic coefficients
\((g_{11},g_{12},g_{44})\).  In compact tensor notation,
\[
  f_{\mathrm{grad}}
  =
  \frac{1}{2}G_{ijkl}
  \frac{\partial P_i}{\partial x_j}
  \frac{\partial P_k}{\partial x_l},
\]
where cubic symmetry reduces \(G_{ijkl}\) to the three fitted coefficients.

Mechanical equilibrium is solved with small strain
\[
  \eps_{ij}
  =
  \frac{1}{2}\left(\frac{\partial u_i}{\partial x_j}
  + \frac{\partial u_j}{\partial x_i}\right),
\]
and the elastic energy is
\[
  f_{\mathrm{el}}
  =
  \frac{1}{2}C_{ijkl}
  \eps^{e}_{ij}\eps^{e}_{kl}.
\]
The elastic strain includes electrostrictive and hidden eigenstrain components:
\[
  \eps^{e}_{ij}
  =
  \eps_{ij}
  - \eps^{Q}_{ij}(\vect{P})
  - \eps^{\mathrm{iso}}_{ij}(\vect{x})
  - \eps^{\mathrm{aniso}}_{ij}(\vect{x}),
\]
with electrostriction represented schematically as
\[
  \eps^{Q}_{ij} = Q_{ijkl}P_kP_l .
\]
The hidden-off baseline has
\[
  \eps^{\mathrm{iso}}_{ij}=0,\qquad
  \eps^{\mathrm{aniso}}_{ij}=0,\qquad
  \lambda_{\mathrm{scr}}=0,\qquad
  h_{\mathrm{flexo}}=0 .
\]

\section{Hidden-Physics Parameterizations}

All hidden-physics candidates were deliberately low-dimensional.  The goal was
not to fit a pixelwise field by BO, but to test physically interpretable
families that could explain structured residuals.

\subsection{Surface Screening / Depolarization Proxy}

The surface-screening field is represented by a localized Gaussian envelope:
\[
  \lambda_{\mathrm{scr}}(\vect{x})
  =
  \lambda_0
  +
  A_{\mathrm{scr}}
  \exp\left[
    -\frac{(x-x_0)^2+(y-y_0)^2}{2\sigma_{xy,\mathrm{scr}}^2}
  \right]
  \exp\left[
    -\frac{(z-z_0)^2}{2\sigma_{z,\mathrm{scr}}^2}
  \right].
\]
In the latest holdout model,
\[
  A_{\mathrm{scr}}=0.08,\qquad
  \sigma_{xy,\mathrm{scr}}=80~\mathrm{nm},\qquad
  \sigma_{z,\mathrm{scr}}=25~\mathrm{nm}.
\]
In the MOOSE input this appears as \texttt{spatial\_screen\_lambda} and enters
the \(P_z\) equation through \texttt{screening\_depol\_z}.  The residual-level
implementation is proportional to
\[
  R^{\mathrm{scr}}_{P_z}
  \propto
  -\frac{1}{2\epsilon_b}
  \lambda_{\mathrm{scr}}(\vect{x})\,\langle P_z\rangle,
\]
with \(\epsilon_b=\texttt{screen\_permitivitty}=0.08854187\).  This term should
be read as an effective depolarization/screening proxy rather than a complete
microscopic charge-transport model.

\subsection{Isotropic Vegard-Like Eigenstrain}

The isotropic Vegard-like field was included as
\[
  \eps^{\mathrm{iso}}_{ij}(\vect{x})
  =
  \eta_{\mathrm{iso}}(\vect{x})\delta_{ij},
\]
where
\[
  \eta_{\mathrm{iso}}(\vect{x})
  =
  \eta_0
  +
  A_{\mathrm{iso}}
  \exp\left[
    -\frac{(x-x_0)^2+(y-y_0)^2}{2\sigma_{xy,\mathrm{iso}}^2}
  \right]
  \exp\left[
    -\frac{(z-z_0)^2}{2\sigma_{z,\mathrm{iso}}^2}
  \right].
\]
In the final anisotropic holdout winner the isotropic Vegard amplitudes are
zero.  Thus the final interpretation is not a uniform volumetric vacancy
expansion/contraction; it is a more directional elastic signature.

\subsection{Low-Rank Anisotropic Eigenstrain}

The anisotropic hidden field is a rank-restricted symmetric eigenstrain tensor:
\[
  \eps^{\mathrm{aniso}}(\vect{x})
  =
  \Phi_A(\vect{x})\,\vect{A},
\]
where
\[
  \Phi_A(\vect{x})
  =
  \exp\left[
    -\frac{(x-x_0)^2+(y-y_0)^2}{2\sigma_{xy,A}^2}
  \right]
  \exp\left[
    -\frac{(z-z_0)^2}{2\sigma_{z,A}^2}
  \right],
\]
and
\[
  \vect{A}
  =
  \begin{pmatrix}
    A_{xx} & A_{xy} & A_{xz}\\
    A_{xy} & A_{yy} & A_{yz}\\
    A_{xz} & A_{yz} & A_{zz}
  \end{pmatrix}.
\]
The MOOSE implementation uses separate \texttt{ComputeEigenstrain} blocks for
\(xx,yy,zz,xy,xz,yz\) and adds all of them to
\texttt{eigenstrain\_names} in the small-strain material.

The best final candidate has only the \(yz\) component active:
\[
  A_{yz}=7.5\times 10^{-4},
  \qquad
  A_{xx}=A_{yy}=A_{zz}=A_{xy}=A_{xz}=0,
\]
with
\[
  \sigma_{xy,A}=120~\mathrm{nm},\qquad
  \sigma_{z,A}=35~\mathrm{nm}.
\]
This is a localized shear-like eigenstrain.  A physically cautious
interpretation is that the residual field is consistent with asymmetric
defect/dislocation-mediated elastic accommodation, rather than with a purely
isotropic vacancy Vegard strain.

\subsection{Phenomenological Flexo Proxy}

A full flexoelectric model would include Lifshitz-type couplings between
polarization gradients and strain gradients, for example terms of the form
\[
  f_{\mathrm{flexo}}
  \sim
  \mu_{ijkl}
  \left(
    P_i\frac{\partial \eps_{jk}}{\partial x_l}
    -
    \eps_{jk}\frac{\partial P_i}{\partial x_l}
  \right),
\]
which generally requires additional carefully verified kernels/materials.
The present campaign did not implement that rigorous tensorial flexoelectric
energy.  Instead, it tested a phenomenological internal-field proxy acting on
\(P_z\):
\[
  h_{\mathrm{flexo}}(\vect{x})
  =
  \left[
    a_0 + a_1\frac{z-z_0}{\sigma_{z,f}}
  \right]
  \exp\left[
    -\frac{(x-x_0)^2+(y-y_0)^2}{2\sigma_{xy,f}^2}
  \right]
  \exp\left[
    -\frac{(z-z_0)^2}{2\sigma_{z,f}^2}
  \right].
\]
Equivalently, one may view the proxy as an energy contribution
\[
  \mathcal{F}_{\mathrm{proxy}}
  =
  -\int_{\Omega} h_{\mathrm{flexo}}(\vect{x})P_z(\vect{x})\,\dd\Omega ,
\]
whose variational derivative drives the \(P_z\) equation.  In MOOSE this is
the \texttt{flexo\_proxy\_pz} body force with function
\texttt{flexo\_proxy\_drive}.  The best flexo-proxy competitor used
\[
  a_1=-10^{-3},\qquad a_0=0,
\]
and reached a full-holdout nMSE of \(0.0687001\), slightly worse than the
anisotropic shear-like eigenstrain.

\section{Candidate Ranking and Quantitative Results}

The full 30-condition holdout comparison is:
\[
\begin{array}{llll}
\toprule
\mathrm{rank} & \mathrm{candidate} & \mathrm{active~terms} & \mathrm{mean~nMSE}\\
\midrule
1 & \texttt{shear\_yz\_pos} &
  \lambda_{\mathrm{scr}} + A_{yz}\Phi_A & 0.0660709\\
2 & \texttt{flexo\_proxy\_grad\_neg} &
  \lambda_{\mathrm{scr}} + h_{\mathrm{flexo}} & 0.0687001\\
3 & \texttt{screening\_only\_refined\_control} &
  \lambda_{\mathrm{scr}} & 0.0698688\\
4 & \texttt{fixed\_g\_hidden\_off} &
  \mathrm{none} & 0.0977860\\
\bottomrule
\end{array}
\]
The best anisotropic candidate reduces nMSE by
\[
  \frac{0.0977860-0.0660709}{0.0977860}\times 100
  =
  32.43\%
\]
relative to the hidden-off fixed-\(g\) model.  It improves all
\(30/30\) holdout conditions relative to hidden-off, \(23/30\) relative to the
screening-only control, and \(21/30\) relative to the flexo-proxy competitor.

\section{Scientific Interpretation}

The important point is not only that the residual decreases, but that the
residual decreases after each candidate is run through the MOOSE phase-field
solver.  Thus the optimization is physics-constrained in the practical sense:
the learned correction must survive coupled LGD, electrostatic, and mechanical
relaxation.  The loss does not currently include an explicit free-energy
penalty; free energy is used diagnostically.

The current best hidden physics is therefore best described as:
\[
  \boxed{
    \mathrm{localized~screening}
    +
    \mathrm{localized~anisotropic~} \eps_{yz}^{*}
    \mathrm{~eigenstrain}
  }.
\]
This is consistent with a sample containing defects/dislocations that produce
directional elastic accommodation near the switched/nucleated region.  It is
not, by itself, proof of a unique microscopic defect population.  The
anisotropic eigenstrain should be presented as an inferred effective field
whose symmetry and localization provide a physics cue for follow-up
experiments or higher-fidelity defect models.

\section{Regenerated Figure Files}

The latest manuscript figures were regenerated by
\texttt{paper\_figures/make\_latest\_campaign\_fig3\_fig4\_story\_panels.py}.
The requested copies are:
\begin{itemize}[leftmargin=*]
  \item \texttt{paper\_figures/needs\_regeneration\_for the\_latest\_campaign/Figure3\_hidden\_physics\_discovery.pdf}
  \item \texttt{paper\_figures/needs\_regeneration\_for the\_latest\_campaign/fig4\_holdout\_validation\_redesigned.pdf}
\end{itemize}
Figure 3a now includes the hidden-off nMSE and intentionally omits error bars.
The two holdout residual atlas panels in Figure 3 use the same center-crop
z-normalized residual color scale,
\[
  R_c \in [-1.2624590595,\;1.2624590595],
\]
so the apparent residual reduction is visually comparable between hidden-off
and hidden-on panels.

\end{document}
